\documentclass[11pt]{article}
\usepackage[margin=1in]{geometry}
\usepackage{amsmath,amssymb,amsthm}
\usepackage{booktabs}
\usepackage{natbib}
\usepackage{hyperref}
\newtheorem{proposition}{Proposition}
\newtheorem{lemma}{Lemma}
\title{Standardization Scope and the Direction of Innovation}
\author{}
\date{}
\begin{document}
\maketitle
\begin{abstract}
We study a regulator that chooses the scope of a common technical standard before competing firms allocate a fixed R\&D capacity between common-layer and proprietary innovation. A broader standard makes common-layer improvements more transferable across rival products. This raises the social return to common innovation but lowers its private return because a firm's common R\&D strengthens its competitor. Firms therefore redirect R\&D toward proprietary innovation as standardization expands, whereas a social planner reallocates R\&D in the opposite direction. In a differentiated Bertrand model, complete standardization is optimal for every fixed positive common-R\&D allocation and under the first best. With decentralized endogenous R\&D allocation, however, sufficiently strong product-market rivalry yields a unique interior standardization scope. Selective standardization is therefore a second-best response to an innovation-composition distortion rather than evidence that interoperability itself is socially harmful.
\end{abstract}

\noindent\textbf{Keywords:} technical standards; interoperability; R\&D allocation; innovation direction; product-market competition; second-best policy.\\
\noindent\textbf{JEL:} L15, L13, O31, O38.

\section{Introduction}

Technical standards can increase interoperability and enlarge the set of products that can use a common technological improvement. But firms do not passively accept the competitive environment created by a standard. When a broader standard makes an improvement in the common layer more useful to rival products, a firm may respond by moving scarce innovative effort toward technologies that remain proprietary. The policy question is therefore not only how much interoperability a standard creates directly, but also how standard scope changes the composition of innovation.

We study that feedback in a deliberately stripped-down industrial-organization model. A regulator first chooses the scope of a mandatory common technical standard. Two firms then allocate a fixed R\&D capacity between common-layer and proprietary innovation, after which they compete in differentiated Bertrand prices. Common-layer R\&D improves the innovating firm's product and, to an extent determined by the mandated scope, the rival's standard-compliant product. Proprietary R\&D benefits only the innovating firm. Scope therefore changes the relative appropriability of two uses of a fixed innovation budget without directly changing product substitutability.

The first result is an endogenous portfolio response. In the quadratic baseline, broader scope strictly reduces private common-layer R\&D and increases proprietary R\&D. The directional result is more general but weaker: for any differentiable increasing strictly concave innovation technology, the private common-layer allocation is globally nonincreasing in scope, while a coordinated symmetric allocation is globally nondecreasing. If two compared scope choices both induce interior optima, the corresponding order inequalities are strict. We do not claim a pointwise derivative sign for arbitrary general technologies. This distinction matters because corner solutions can generate flat intervals and differentiability of the innovation technology alone does not justify differentiating the optimizer with respect to policy.

The main policy result is deliberately baseline-specific. If the R\&D composition is held fixed at any positive symmetric common-layer allocation, welfare is increasing in scope and complete standardization is optimal. Complete standardization is also optimal when the regulator chooses scope together with the symmetric R\&D composition while differentiated Bertrand pricing remains decentralized. That benchmark is constrained and is not an unconstrained first best. With decentralized endogenous R\&D composition, broader scope induces firms to retreat from common innovation. In the quadratic baseline, sufficiently strong product-market rivalry makes the reduced policy objective strictly concave with a unique interior optimum. Figure~\ref{fig:regime-map} summarizes the exact rivalry threshold. Selective standardization is therefore a scope-only second-best response to an innovation-composition distortion, not a claim that interoperability is intrinsically socially costly.

The contribution is deliberately result-level rather than thematic. The closest technical-standards comparison is \citet{Llanes2024}, who studies incentives to develop technologies for inclusion in a standard together with patent pools, price caps, and cooperative R\&D. Related work shows that compatibility and standard-setting institutions affect innovation incentives \citep{BondSmith2019,MaruyamaZennyo2017,HeywoodWangYe2022}, while endogenous product design can overturn compatibility conclusions \citep{Boom2001} and government standardization can interact with product characteristics \citep{Ruiz2004}. Our policy object is different: a regulator chooses a continuous standardization scope that changes cross-product transferability before two firms reallocate a fixed common-versus-proprietary R\&D portfolio. \citet{AcemogluGanciaZilibotti2012} already characterize constrained optimal standardization in a dynamic growth environment, so a generic second-best case for limiting standardization is not novel here. Likewise, \citet{BryanLemus2017} develop a theory of the direction of innovation with scarce research resources and underappropriation; endogenous innovation direction itself is therefore not our novelty. The paper's contribution is the combination of regulator-chosen scope, endogenous common-versus-proprietary portfolio reallocation, differentiated-product rivalry, and the exact complete-to-selective policy reversal in the quadratic baseline.

The mechanism is intended for settings with a meaningful common interface or standardized technological layer and a distinct proprietary layer. Digital protocols and APIs and modular industrial interfaces are natural interpretations. Evidence from modular product development is consistent with innovative attention shifting toward internal interfaces and modules when external interfaces are standardized \citep{ChenLiu2005}. Recent firm-level evidence also shows that firms technologically closer to a newly adopted standard increase R\&D more strongly in competitive markets and that standardization can create a common technological basis that facilitates later catch-up \citep{BergeaudSchmidtZago2026}. We use these findings only to motivate the interaction among standardization, competition, and innovation; they are not causal validation of the portfolio mechanism modeled here.

The rest of the paper proceeds as follows. Section~2 defines the game and the all-history price-continuation restriction. Section~3 characterizes private R\&D allocation and derives the selective-standardization threshold. Section~4 derives welfare and compares the decentralized policy problem with fixed-symmetric and coordinated-symmetric benchmarks without claiming an unconstrained first best. Section~5 states only the robustness results authorized by the theory freeze. Section~6 positions the contribution relative to compatibility, standards-and-innovation, innovation-direction, and R\&D-allocation literatures. The Appendix provides global price-continuation, corner/KKT, threshold, and verification details.

\section{Model}

There are two symmetric firms, indexed by $i\in\{1,2\}$, a benevolent regulator, and a representative consumer. The game has three stages. First, the regulator chooses the scope of a mandatory common technical standard, $b\in[0,1]$. Second, firms simultaneously allocate a fixed R\&D capacity $E>0$ between common-layer innovation $x_i$ and proprietary innovation $z_i$, subject to
\begin{equation}
 x_i+z_i=E,\qquad x_i,z_i\geq 0.
\end{equation}
Third, firms simultaneously choose nonnegative prices $p_i\geq0$.

Innovation in either layer raises product quality through
\begin{equation}
 g(r)=r-\frac{\kappa}{2}r^2,
\end{equation}
with $\kappa>0$. Define $y\equiv \kappa E$ and impose
\begin{equation}
 \frac12<y<1.
 \label{eq:yregion}
\end{equation}
Hence $g'(r)=1-\kappa r>0$ on $[0,E]$ and $g$ is strictly concave.

The quality intercept of firm $i$ is
\begin{equation}
 A_i=a+\theta\left[g(x_i)+b g(x_j)+g(E-x_i)\right],\qquad j\neq i,
 \label{eq:quality}
\end{equation}
where $a>0$ and $\theta>0$. The term $g(x_i)$ is the direct benefit of the firm's own common-layer R\&D. The term $b g(x_j)$ captures the fraction of the rival's common-layer improvement that becomes usable in firm $i$'s standard-compliant product as the mandated common scope expands. Proprietary R\&D, $g(E-x_i)$, benefits only the innovating firm. Thus $b$ changes the transferability of common innovation; it does not directly alter product substitutability.

The representative consumer has quasi-linear utility over the two products,
\begin{equation}
 U(q_1,q_2)=A_1q_1+A_2q_2-\frac12(q_1^2+q_2^2)-\rho q_1q_2,
 \label{eq:utility}
\end{equation}
with $q_i\geq0$ and $0<\rho<1$. Production has zero marginal cost. In the region where both products are consumed, demand is
\begin{equation}
 q_i=\frac{A_i-p_i-\rho(A_j-p_j)}{1-\rho^2}.
 \label{eq:demand}
\end{equation}

The baseline equilibrium concept is subgame-perfect Nash equilibrium. Because the R\&D and policy choices are upstream of price competition, we require a well-defined price continuation at every feasible upstream history. Let
\begin{equation}
 A_{\min}=a+\theta E\left(1-\frac y2\right),\qquad
 A_{\max}=a+\theta E\left(2-\frac{3y}{4}\right).
\end{equation}
We impose the sufficient regularity condition
\begin{equation}
 \frac{A_{\max}}{A_{\min}}
 <
 \frac{2}{\rho(3-\rho^2)}.
 \tag{R}\label{eq:R}
\end{equation}
Condition~\eqref{eq:R} guarantees that the unique Bertrand equilibrium keeps both products active for every $b\in[0,1]$ and every $(x_1,x_2)\in[0,E]^2$, and that no nonnegative unilateral price deviation can profitably foreclose the rival. It is a global-continuation regularity condition rather than a condition needed for the economic mechanism itself.

It is useful to define
\begin{equation}
 \nu\equiv\frac{\rho}{2-\rho^2}.
 \label{eq:nu}
\end{equation}
Since $0<\rho<1$, we have $0<\nu<1$. The baseline interior R\&D region additionally imposes
\begin{equation}
 0<\nu<y.
 \label{eq:nuy}
\end{equation}
As shown below, $\nu$ measures the competitive harm from improving the rival relative to the private return from improving one's own product.

\section{Equilibrium and the Direction of Innovation}

We solve the game backward. Under condition~\eqref{eq:R}, the price subgame has a unique active-product equilibrium at every feasible upstream history. The first-order conditions are
\begin{equation}
 2p_i-\rho p_j=A_i-\rho A_j,
\end{equation}
which yield
\begin{equation}
 p_i^*=\frac{(2-\rho^2)A_i-\rho A_j}{4-\rho^2}.
 \label{eq:price}
\end{equation}
Demand satisfies $q_i^*=p_i^*/(1-\rho^2)$, so equilibrium profit is
\begin{equation}
 \pi_i^*=\frac{\left[(2-\rho^2)A_i-\rho A_j\right]^2}
 {(4-\rho^2)^2(1-\rho^2)}.
 \label{eq:profit}
\end{equation}
Condition~\eqref{eq:R} guarantees that the bracketed term is positive for all feasible histories. Hence maximizing profit with respect to $x_i$ is equivalent to maximizing the linear quality index $(2-\rho^2)A_i-\rho A_j$.

Using \eqref{eq:quality} and dividing by the positive coefficient $2-\rho^2$, firm $i$ solves
\begin{equation}
 \max_{x_i\in[0,E]}
 \left(1-\nu b\right)g(x_i)+g(E-x_i),
 \label{eq:rdproblem}
\end{equation}
where $\nu$ is defined in \eqref{eq:nu}. The rival's R\&D choice drops out of firm $i$'s best response. Economically, one unit of common-layer R\&D improves the firm's own product but, through the standard, also improves the rival's product. The term $\nu b$ is the private competitive penalty attached to that spillover.

The objective in \eqref{eq:rdproblem} is strictly concave. Under \eqref{eq:nuy}, the unique optimum is interior for every $b\in[0,1]$ and equals
\begin{equation}
 x^F(b)=\frac{y-\nu b}{\kappa(2-\nu b)}.
 \label{eq:xprivate}
\end{equation}
Because each firm's best response is independent of the rival's allocation, the R\&D subgame has a unique symmetric equilibrium and no asymmetric R\&D equilibrium in the frozen parameter region.

\begin{proposition}[Private R\&D reallocation]\label{prop:private}
Suppose \eqref{eq:yregion}, \eqref{eq:nuy}, and \eqref{eq:R} hold. The unique private R\&D allocation satisfies
\begin{equation}
 \frac{dx^F}{db}
 =-\frac{\nu(2-y)}{\kappa(2-\nu b)^2}<0,
 \qquad
 \frac{dz^F}{db}>0.
\end{equation}
Thus a broader common standard redirects a fixed R\&D budget from common-layer toward proprietary innovation.
\end{proposition}

The sign does not rely on the quadratic specification. For any increasing strictly concave $g$, the private first-order condition is
\begin{equation}
 (1-\nu b)g'(x^F)=g'(E-x^F),
 \label{eq:privategeneral}
\end{equation}
and implicit differentiation gives $dx^F/db<0$.

For comparison, consider a planner that, conditional on a given scope $b$, chooses the symmetric R\&D composition to maximize symmetric product quality and therefore welfare. The planner solves
\begin{equation}
 \max_{x\in[0,E]} (1+b)g(x)+g(E-x),
\end{equation}
with first-order condition
\begin{equation}
 (1+b)g'(x^S)=g'(E-x^S).
 \label{eq:socialgeneral}
\end{equation}
Under the quadratic technology,
\begin{equation}
 x^S(b)=\frac{y+b}{\kappa(2+b)}.
 \label{eq:xsocial}
\end{equation}

\begin{proposition}[Directional private--social wedge]\label{prop:wedge}
For any increasing strictly concave $g$, a broader standard shifts the private and social R\&D allocations in opposite directions:
\begin{equation}
 \frac{dx^F}{db}<0<\frac{dx^S}{db}.
\end{equation}
Moreover, for $b>0$,
\begin{equation}
 x^F(b)<\frac{E}{2}<x^S(b).
\end{equation}
\end{proposition}

Proposition~\ref{prop:wedge} is the mechanism, not the main policy result. A standard that widens the use of common innovation makes such innovation less privately attractive because it strengthens competitors, while making it more socially attractive because the same technological improvement benefits a larger set of products.

\subsection{The regulator's scope choice}

We first hold the R\&D composition fixed at an arbitrary $\bar x>0$. Symmetric quality is then
\begin{equation}
 A^{BT}(b)=a+\theta\left[(1+b)g(\bar x)+g(E-\bar x)\right].
\end{equation}
Since $\partial A^{BT}/\partial b=\theta g(\bar x)>0$, welfare is increasing in $b$ for every fixed allocation with positive common R\&D.

\begin{proposition}[Fixed-allocation benchmark]\label{prop:benchmark}
For every fixed R\&D allocation with $\bar x>0$, the regulator chooses complete standardization:
\begin{equation}
 b^{BT}=1.
\end{equation}
If $\bar x=0$, the regulator is indifferent over $b$.
\end{proposition}

Now substitute the private R\&D equilibrium \eqref{eq:xprivate}. Since symmetric welfare is strictly increasing in symmetric quality, the policy problem is equivalent to maximizing
\begin{equation}
 F(b)=\kappa\left[(1+b)g(x^F(b))+g(E-x^F(b))\right].
 \label{eq:F}
\end{equation}
Direct calculation gives
\begin{equation}
 F''(b)=
 -\frac{\nu(2-y)^2(2b\nu^2+b\nu+2\nu+4)}{(2-b\nu)^4}<0,
 \label{eq:Fpp}
\end{equation}
and
\begin{equation}
 F'(0)=\frac{y(4-y)}{8}>0.
 \label{eq:Fp0}
\end{equation}
Thus the regulator never chooses zero standardization in the baseline region. At $b=1$, define
\begin{align}
 H(y,\nu)={}&\nu^3+2\nu^2y^2-8\nu^2y+2\nu^2 \\
 &+\nu y^2-4\nu y+16\nu+2y^2-8y.
 \label{eq:H}
\end{align}
The sign of $F'(1)$ is equivalent to the sign of $H$: $F'(1)<0$ if and only if $H(y,\nu)>0$. Furthermore,
\begin{equation}
 H(y,0)=2y(y-4)<0,
 \qquad
 H(y,y)=2y(y-2)^2(y+1)>0,
\end{equation}
and $H_\nu>0$ on $0<\nu<y<1$. Hence for every $y\in(1/2,1)$ there is a unique threshold $\bar\nu(y)\in(0,y)$ satisfying $H(y,\bar\nu(y))=0$.

\begin{proposition}[Selective standardization]\label{prop:selective}
Suppose \eqref{eq:yregion}, \eqref{eq:nuy}, and \eqref{eq:R} hold. For each $y\in(1/2,1)$ there exists a unique $\bar\nu(y)\in(0,y)$ such that
\begin{equation}
 \nu\leq\bar\nu(y) \quad\Longrightarrow\quad b^{FULL}=1,
\end{equation}
while
\begin{equation}
 \bar\nu(y)<\nu<y \quad\Longrightarrow\quad 0<b^{FULL}<1.
\end{equation}
In the latter region the interior policy is unique.
\end{proposition}

Because $\nu=\rho/(2-\rho^2)$ is increasing in $\rho$, stronger product-market substitutability makes selective standardization more likely. Rivalry increases the private cost of transferring the benefit of common R\&D to a competitor. The regulator responds not because the direct value of standardization has become negative, but because complete standardization induces too large a retreat from common innovation.

For illustration, let $\kappa=1$, $E=0.7$, $\nu=0.5$, $\rho=\sqrt{3}-1$, $a=10$, and $\theta=0.2$. The verified numerical solution is $b^{FULL}\simeq0.6878$, whereas the fixed-allocation benchmark selects $b=1$.

\section{Welfare}
Stage 10 writing placeholder. Theory is frozen.

\section{Approved Robustness and Interpretation}

This section records only robustness and interpretation statements authorized by the active theory freeze. It does not extend the exact selective-standardization theorem beyond the quadratic baseline.

\subsection{General concave innovation technology}

Let $g$ be differentiable, increasing, and strictly concave on $[0,E]$. As Proposition~\ref{prop:wedge} states, for any $b_2>b_1$ the private common-layer allocation satisfies
\begin{equation}
 x^F(b_2)\leq x^F(b_1),
\end{equation}
while the coordinated symmetric allocation satisfies
\begin{equation}
 x^S(b_2)\geq x^S(b_1).
\end{equation}
These global order statements include corner solutions. If both compared private optima are interior, the first inequality is strict; if both compared coordinated optima are interior, the second inequality is strict. The interior first-order conditions are
\begin{equation}
 (1-\nu b)g'(x^F)=g'(E-x^F),
 \qquad
 (1+b)g'(x^S)=g'(E-x^S),
\end{equation}
but they are used only to establish strict ordering across two policy values. We do not assert a pointwise sign for $dx^F/db$ or $dx^S/db$ for arbitrary general technologies.

The exact polynomial threshold and global uniqueness in Proposition~\ref{prop:selective} rely on the quadratic baseline. We make no generic claim that sufficiently small $C^2$ perturbations preserve strict concavity or uniqueness of the regulator's reduced objective: induced policy curvature can depend on higher-order features of the innovation technology. The general-technology robustness result is therefore the directional allocation ordering, not a general policy theorem.

\subsection{Incomplete transferability}

Suppose common-layer transferability is $t=\lambda b$ with $\lambda\in(0,1]$. Replacing $b$ by $t$ in the relevant quality spillover gives the interior conditions
\begin{equation}
 (1-\nu t)g'(x^F)=g'(E-x^F),
 \qquad
 (1+t)g'(x^S)=g'(E-x^S).
\end{equation}
Thus the same private-versus-coordinated relative-return wedge is indexed by effective transferability. This reparameterization is the approved robustness statement. We do not infer an additional global policy-uniqueness or persistence theorem from it.

\subsection{Endogenous total R\&D}

The fixed-capacity assumption isolates the composition margin. Consider only an extension in which firms choose common effort $x$ and proprietary effort $z$ and face the same marginal total-capacity cost through a term such as $C(x+z)$. At an interior optimum, that common marginal cost appears in both activity first-order conditions and cancels when the two conditions are compared. The resulting relative condition is
\begin{equation}
 (1-\nu b)g'(x)=g'(z).
 \label{eq:endogenousErelative}
\end{equation}
This condition shows that broader scope still lowers the private relative return to common effort. It does not determine the response of total R\&D, the global response of common effort, or the optimal standardization scope. In particular, we do not claim that a steep capacity cost makes the selective-standardization policy theorem persist by continuity. Those questions require a separate model and are outside the frozen result set.

\subsection{Institutional interpretation}

The variable $b$ is the mandatory scope of a common technological layer or interface. It is not a generic compatibility index and does not mean that all firm knowledge becomes legally disclosed. The reduced-form spillover in \eqref{eq:quality} represents greater cross-product usability of improvements made within the standardized layer.

Two application families illustrate the mechanism without changing the theory. In digital systems, $x$ may represent protocol-, API-, or interface-facing innovation and $z$ proprietary implementation or service innovation. In modular manufacturing, $x$ may represent innovation in standardized external interfaces and $z$ innovation in internal modules or proprietary subsystems. Case evidence from Taiwan's machine-tool industry reports that external interface standards shape product-innovation choices and that innovative effort can concentrate on internal interfaces and modules \citep{ChenLiu2005}. We treat this as suggestive institutional consistency, not causal validation.

The empirical implications should respect the theorem's scope. In the quadratic interior baseline, broader standard scope strictly lowers private common-layer R\&D, and the response is stronger when the competitive penalty $\nu$ is larger. Under a general differentiable increasing strictly concave technology, the corresponding prediction is an order statement: private common-layer R\&D is nonincreasing in scope and may remain flat over ranges where a corner binds. In the quadratic baseline, because $\nu$ rises with downstream substitutability $\rho$, stronger rivalry moves the economy toward the selective-standardization region characterized in Figure~\ref{fig:regime-map}. These predictions concern innovation composition rather than total R\&D expenditure. The model does not imply that the realized transferred improvement $b g(x^F(b))$ must itself rise monotonically with scope.

\section{Related Literature and Contribution Boundary}

The paper sits at the intersection of three literatures: compatibility and endogenous product design, innovation incentives under standards and compatibility, and the allocation of innovative effort across technological activities.

First, compatibility need not be welfare improving once product design is endogenous. \citet{Boom2001} studies a three-stage compatibility--product-specification--pricing game and shows that endogenous product selection can overturn the welfare gains from compatibility. \citet{Ruiz2004} studies international standardization policy with differentiated products and shows that endogenous product characteristics affect the consequences of recognizing foreign standards. These papers establish that firms can redesign products in response to compatibility policy. Our adjustment margin is different: product substitutability is fixed, while firms reallocate a scarce R\&D budget between a common technological layer and a proprietary layer.

Second, a substantial literature links compatibility to innovation incentives. \citet{BondSmith2019} studies compatibility, endogenous R\&D, and commercialization in markets with network effects. \citet{MaruyamaZennyo2017} analyzes process innovation and endogenous application compatibility, showing that improved R\&D efficiency can induce lower compatibility and reduce welfare. \citet{HeywoodWangYe2022} compares R\&D rivalry and cooperation when compatibility is endogenous. \citet{Llanes2024} studies incentives to develop technologies for inclusion in a technical standard and shows that standardization can generate either insufficient or excessive innovation. In a dynamic general-equilibrium setting, \citet{AcemogluGanciaZilibotti2012} analyzes the competing roles of innovation and standardization and characterizes an optimal speed of standardization. These contributions make clear that neither ``compatibility affects innovation'' nor ``partial standardization may be optimal'' is novel by itself.

Third, the allocation of innovative effort across alternative activities is also established. \citet{KretschmerReitzig2013} model firms' R\&D diversification in emerging standard settings, where firms trade off cooperation on a joint standard against own activities. \citet{daSilva2024} studies scarce managerial resources allocated across internal and external R\&D and shows how spillovers and substitutability across research approaches affect the allocation. Evidence on modular product innovation likewise suggests that external interface standards can redirect innovative attention toward internal interfaces and modules \citep{ChenLiu2005}.

The contribution of the present paper is therefore deliberately narrower than any of those broad statements. We study a regulator that chooses a continuous standardization scope before firms allocate a fixed innovation capacity. The same scope parameter enters private and social R\&D returns with opposite strategic signs: greater transferability lowers the private relative return to common innovation because it strengthens rivals, while increasing its social relative return because the innovation serves more products. This produces a full-game policy reversal. Holding R\&D composition fixed, complete standardization is optimal; complete standardization is also first-best. Yet once firms endogenously redirect a fixed R\&D portfolio, sufficiently strong product-market rivalry makes the regulator choose a unique interior scope. The novel object is this fixed-allocation/full-standardization to endogenous-portfolio/selective-standardization comparison and its threshold characterization, not the component mechanisms considered separately.

\section{Conclusion}

A broader technical standard can improve interoperability and diffuse common technological improvements across products. But it can also change what firms choose to innovate. This paper isolates that second effect by allowing competing firms to allocate a fixed R\&D capacity between common-layer and proprietary innovation after the regulator chooses the common scope.

The private and social responses move in opposite directions. Firms reduce common-layer R\&D as the standard expands because their improvements increasingly benefit a rival, whereas a social planner increases common-layer R\&D because the same transferability broadens the social return to innovation. This innovation-composition distortion feeds back into the regulator's policy problem. Complete standardization is optimal for every fixed positive common-R\&D allocation and under the first best, but sufficiently strong product-market rivalry yields a unique interior standardization scope when R\&D allocation is decentralized.

The policy interpretation is explicitly second-best. Selective standardization is not desirable because interoperability is intrinsically costly in the model. It is desirable because standardization is the regulator's only instrument and broader scope weakens private incentives to invest in the common layer. If another policy instrument could restore the private return to common innovation, the case for restricting the standard could disappear. Evaluating such instruments is beyond the frozen model.

The broader implication is that the scope of a technical standard and the direction of innovation should be analyzed jointly. A regulator that evaluates standardization while holding the innovation portfolio fixed may favor complete standardization even in an environment where the equilibrium response of firms makes a narrower common scope optimal. Future empirical work could test this mechanism by tracing R\&D or patent portfolios before and after changes in the scope of common technical interfaces, with particular attention to variation in downstream product substitutability.

\appendix
\section{Proof Details and Verification Notes}

\subsection{Global price continuation}
For fixed qualities $(A_i,A_j)$, the active-region profit of firm $i$ is
\begin{equation}
 \pi_i(p_i;p_j)=\frac{p_i\left[A_i-p_i-\rho(A_j-p_j)\right]}{1-\rho^2},
\end{equation}
which is strictly concave in $p_i$. Solving the two first-order conditions gives \eqref{eq:price}. Condition~\eqref{eq:R} implies positive candidate prices and quantities at every feasible upstream history and makes the price required to foreclose the rival negative. Hence no feasible nonnegative unilateral deviation can cross into a profitable rival-foreclosure region. Deviations that make own demand zero yield zero profit and are dominated by the positive equilibrium profit. Thus the active candidate is the unique global Nash equilibrium of every price subgame.

\subsection{Private and social R\&D comparative statics}
For a general increasing strictly concave $g$, define
\begin{equation}
 \Phi_F(x,b)=(1-\nu b)g'(x)-g'(E-x).
\end{equation}
At an interior optimum, $\Phi_F=0$ and
\begin{equation}
 \Phi_{F,x}=(1-\nu b)g''(x)+g''(E-x)<0,
 \qquad
 \Phi_{F,b}=-\nu g'(x)<0.
\end{equation}
Therefore $dx^F/db=-\Phi_{F,b}/\Phi_{F,x}<0$. Similarly, for
\begin{equation}
 \Phi_S(x,b)=(1+b)g'(x)-g'(E-x),
\end{equation}
we have $\Phi_{S,b}=g'(x)>0$ and $\Phi_{S,x}<0$, so $dx^S/db>0$.

\subsection{Threshold uniqueness}
Differentiating \eqref{eq:H} with respect to $\nu$ gives
\begin{equation}
 H_\nu=3\nu^2+4\nu(y^2-4y+1)+y^2-4y+16.
\end{equation}
For $0<\nu<y<1$,
\begin{equation}
 H_\nu>4(y^2-4y+1)+(y^2-4y+16)=5(y-2)^2>0.
\end{equation}
Together with $H(y,0)<0<H(y,y)$, continuity implies a unique root $\bar\nu(y)\in(0,y)$.

\subsection{Exact consumer surplus}
At a symmetric equilibrium, $q_1=q_2=q$. Gross utility is
\begin{equation}
 2Aq-(1+\rho)q^2.
\end{equation}
Net consumer surplus subtracts expenditure $2pq$. Using \eqref{eq:symmetricpq} yields \eqref{eq:CS}. Adding equilibrium profits yields \eqref{eq:W}.

\subsection{Computational verification}
The repository contains symbolic identity checks in \texttt{scripts/symbolic\_verify.py}, an independent numerical evaluator in \texttt{scripts/numerical\_verify.py}, and regression tests in \texttt{tests/test\_freeze\_regressions.py}. These artifacts implement the frozen theory record \texttt{SSDI-THEORY-FREEZE-2026-09-06-v1}; they are verification aids rather than substitutes for the analytic arguments above.

\bibliographystyle{apalike}
\bibliography{../references/references}
\end{document}
